\scanpage{28}
\mspage{15}

% top two lines
\[
\begin{aligned}
\text{donc}\quad\mathrm{rg}\,N(u_*)=2&\Rightarrow
 \mathrm{rg}(u_i,u_{i+1}+u_{i+2},u_{i+3})=1\\
&&\forall i
\end{aligned}
\]

% below top implication
Donc $a)\Rightarrow b)$. Un calcul \unsure{fait|facile} plus haut montre

% below Un calcul
\[
b)\Longleftrightarrow c)
\]

% next two lines
et on voit \unsure{tout} de suite que \unsure{c'est|c)|ceci} \unsure{équivt|équivaut} à $d)$.

% centre
\unsure{Or|On} considérons les $(\lambda_i)$ \uline{périodiques} de période $3$, tels que

% above main relation
\[
\deleted{\lambda_0u_0+\lambda_1u_1+\lambda_2}
\]

% centre, relation and lower annotations
\[
\lambda_i u_{i-1}\deleted{\unsure{+}}-\lambda_{i+1}
(\underbrace{u_i+u_{i+1}}_{-u_{i-1}})+
\underset{\substack{\Vert\\\lambda_{i-1}}}{\lambda_{i+2}}
(\underset{\substack{\Vert\\u_{i-1}}}{u_{i+2}})=0\qquad\forall i
\]

% centre, reduction
\[
\text{i.e.}\quad
(\lambda_i+\lambda_{i-1}\deleted{\unsure{u_{i-1}|\lambda_{i-1}}}+\lambda_{i+1})u_{i-1}=0
\]

% final two lines
Ils \unsure{font|forment} un \unsure{vectoriel} de dim $2$ ! Donc $d)\Rightarrow a)$

% archive sheet number: 34
